% SOURCE docs/release/2026-10-29/source/docs/writing-repair/2026-10-29/src/cfe939aa8760d37c.tex CHAR 77
\chapter{Persistent audit information and finite calendar
uncertainty}\label{persistent-audit-information-and-finite-calendar-uncertainty}}

\label{exec13:persistent-region}

A bank can increase its account count without acquiring an independent
observation of a common calendar shock. It can also add monthly
observations without redrawing a persistent identity audit. These two
limitations affect different parts of funding uncertainty. A useful
analysis must first state the joint observation law, then determine
which uncertainty can be bounded from its actual information.

The construction below is a project derivation for the additive
experiment. It gives a finite-sample confidence set when variance and
persistence bounds are supplied, together with an information
counterexample to a uniform precision claim when persistence is
unrestricted. The exact multiplier moments, covariance references and
numerical illustrations are reproducible project calculations; no new
random evaluation sample or empirical qualification is used.
\citep{exec13_persistent_derivation,exec13_persistent_examples}

\hypertarget{the-observation-law-and-the-banking-target}{%
\section{The observation law and the banking
target}\label{the-observation-law-and-the-banking-target}}

Retain the previously studied independent persistent pairs, common
calendar and separate gold experiment. For pair \(i=1,\ldots,N\) and
date \(t=1,\ldots,T\), the signed focal and companion responses are

\[
 F_{it}=\beta+U_i+C_t+E_{it},\qquad
 D_{it}=\delta+V_i+C_t+J_{it}.
\]

The pair effects \(U_i,V_i\) are mutually independent centered Gaussians
with variance \(v_p\). The cell errors \(E_{it},J_{it}\) are mutually
independent centered Gaussians with variance \(v_e\). The independent
calendar is stationary Gaussian AR(1) with variance \(v_c\), correlation
\(\rho^{|t-s|}\), and \(|\rho|<1\). Candidate order is fixed within a
pair and exogenous. Its sign \(K_i\in\{-1,1\}\) states whether the first
candidate is focal. Randomized observed action signs multiply both
recorded candidate histories; undoing those known signs yields the two
candidate responses \(X_{i1t},X_{i2t}\).

An audit \(R_i\sim\operatorname{Bernoulli}(q)\), with known constant
\(q>0\), selects a whole pair history. On selected pairs the reported
order is \(L_i=K_i\xi_i\), where \(\xi_i\) has probability \(e<1/2\) of
being \(-1\). The error is symmetric, independent of pair outcomes and
audits, and persistent across dates. An independent representative gold
sample contains \(M\) Bernoulli error indicators with the same \(e\).
Unobserved audit labels remain missing; multiplication by \(R_i\) below
means their contribution is zero without imputing an observed label.
Endogenous selection, asymmetric error, shared pair audits and
nonrepresentative gold labels violate this law.

The target is \(1000\beta\) currency units per focal unit under the
previously studied signed-action response convention. This causal
interpretation is stipulated in the generated experiment; multiplying an
arbitrary bank outcome by an action sign would not establish it. Let
\(b=(\beta+\delta)/2\), \(d=(\beta-\delta)/2\), and \(r=1-2e>0\). Then
\(\beta=b+d\).

The two observed summaries are

\[
 \widehat B=\frac1{NT}\sum_{it}\frac{X_{i1t}+X_{i2t}}2,
 \qquad
 \widehat H=\frac1N\sum_i H_i,
 \quad
 H_i=\frac{R_iL_i}{2qT}\sum_t(X_{i1t}-X_{i2t}).
\]

The common calendar cancels from every candidate difference. Since the
same \(K_i\) appears in the difference and the audited order,
\(K_i^2=1\) removes ordering from \(H_i\). Hence

\[
 H_i=W_i\left(d+\frac{U_i-V_i}{2}
                    +\frac1T\sum_t\frac{E_{it}-J_{it}}2\right),
 \qquad W_i=R_i\xi_i/q,
\]

with \(E W_i=r\) and \(E W_i^2=1/q\). The independent-pair and Gaussian
assumptions give

\[
 E\widehat B=b,\qquad E\widehat H=rd,
\]

\[
 V_B=\operatorname{Var}(\widehat B)
   =\frac{v_p}{2N}+\frac{v_e}{2NT}+v_c k_T(\rho),
 \quad
 k_T(\rho)=\frac{T+2\sum_{k=1}^{T-1}(T-k)\rho^k}{T^2},
\]

\[
 \operatorname{Var}(\widehat H)
 =\frac1N\left\{\frac{v_p}{2q}+\frac{v_e}{2qT}
                       +(q^{-1}-r^2)d^2\right\}.
\]

These formulas do not divide by an estimated rank. The first summary is
Gaussian. The second is a persistent-audit mixture, so a Gaussian
approximation to it is an additional step, not an exact consequence of
Gaussian outcomes. The coefficient of \(d^2\) retains the audit
uncertainty that more dates cannot remove. The variances are nonnegative
under \(q\leq1\) and \(0<r\leq1\).

\hypertarget{a-finite-construction-with-explicit-unknown-nuisance-bounds}{%
\section{A finite construction with explicit unknown-nuisance
bounds}\label{a-finite-construction-with-explicit-unknown-nuisance-bounds}}

Suppose the admissible model supplies finite bounds \(v_p\leq\bar v_p\),
\(v_e\leq\bar v_e\), \(v_c\leq\bar v_c\) and
\(|\rho|\leq\bar\rho\leq1\). The values inside these bounds are unknown.
Taking \(\bar\rho=1\) is allowed as a conservative envelope for the
stationary models, although it sacrifices calendar shrinkage. Since the
lag coefficients in \(k_T\) are positive,

\[
 V_B\leq \bar V_B:=\frac{\bar v_p}{2N}
            +\frac{\bar v_e}{2NT}+\bar v_c k_T(\bar\rho).
\]

Choose error budgets \(\alpha_B,\alpha_H,\alpha_e>0\) whose sum is at
most \(\alpha\). The normal interval

\[
 C_B=[\widehat B-z_{1-\alpha_B/2}\sqrt{\bar V_B},
      \widehat B+z_{1-\alpha_B/2}\sqrt{\bar V_B}]
\]

has coverage at least \(1-\alpha_B\), because \(\widehat B-b\) is normal
with variance no larger than \(\bar V_B\). Degenerate zero variance is
handled by the point interval, not by dividing by zero.

Obtain a two-sided exact binomial interval \(C_e\) by inverting the gold
count \(S\sim\operatorname{Binomial}(M,e)\), allocating \(\alpha_e/2\)
to each tail. For an interior count \(s\), its lower endpoint solves
\(P_{e_L}(S\geq s)=\alpha_e/2\) and its upper endpoint solves
\(P_{e_U}(S\leq s)=\alpha_e/2\), with endpoints zero/one for the
boundary counts. Binomial tail monotonicity shows that excluding the
true \(e\) through either tail has probability at most its allocated
error. Intersect with the model's \([0,1/2)\) to obtain
\(C_r=\{1-2e:e\in C_e\cap[0,1/2)\}\). This intersection preserves
coverage under the model, but can yield ranks arbitrarily close to zero.
It does not justify silently imposing a positive lower rank.

Define the observed-data confidence set

\[
 C=\left\{(b,d,r): b\in C_B,\ r\in C_r,
 (\widehat H-rd)^2\leq
 \frac{\bar v_p/(2q)+\bar v_e/(2qT)+(q^{-1}-r^2)d^2}
      {N\alpha_H}\right\}.
\]

At the true \((d,r)\), Chebyshev's inequality bounds the probability of
violating the displayed quadratic inequality by \(\alpha_H\), using the
variance derived above and its nuisance envelope. This follows directly
by applying Markov's inequality to \((\widehat H-rd)^2\): its
expectation is the variance and is no larger than the numerator divided
by \(N\). No independence between the confidence components is needed.
The union bound therefore proves

\[
 P\{(b,d,r)\in C\}\geq1-\alpha,
 \qquad
 P\{1000\beta\in\mathcal S\}\geq1-\alpha,
 \quad \mathcal S=\{1000(b+d):(b,d,r)\in C\}.
\]

This is a finite-sample statement, uniform over the specified nuisance
envelope and the declared observation law. It estimates the gold-error
rank through a confidence set and bounds unknown calendar/pair/cell
variances; it does not treat those unknown values as known oracle
constants. The construction is feasible in the mathematical sense that
all inputs to the inequalities are observed summaries or declared
bounds. A numerically certified implementation must still project over a
continuum of \(r\); a grid that misses a feasible component would not
inherit the theorem. Analytic quadratic roots or interval outer bounds,
including empty, disconnected and unbounded sets, are suitable
implementation routes. A raw point estimator is optional.

If independent pilot data supply simultaneous upper confidence bounds on
the variances or persistence instead of deterministic bounds, their
joint failure probability must receive an additional budget. Coverage
follows on the event that the bounds are valid, with that event's
failure probability added by the union bound. Independence is not
required for this composition argument, although it can simplify the
pilot proof. Supplying bounds calculated from the same data without a
validity argument would merely move the missing uncertainty problem into
another function.

\hypertarget{why-validity-need-not-produce-a-useful-finite-interval}{%
\section{Why validity need not produce a useful finite
interval}\label{why-validity-need-not-produce-a-useful-finite-interval}}

For fixed candidate rank \(r\), multiply the numerator inequality by
\(\alpha_H\) and collect powers of \(d\). The leading quadratic
coefficient is

\[
 a(r)=r^2(\alpha_H+N^{-1})-(qN)^{-1}.
\]

If \(a(r)<0\), sufficiently large positive and negative \(d\) satisfy
the inequality. If \(a(r)=0\), the inequality is linear or constant and
can produce a half-line, the whole line or an empty set. If \(a(r)>0\),
it gives a bounded interval or an empty set at that rank. A sufficient
uniform condition for bounded numerator projection is that every
feasible rank exceed

\[
 r_{\min}^2>\frac1{q(N\alpha_H+1)}.
\]

Merely checking a positive leading coefficient at the estimated rank is
insufficient when smaller ranks remain in \(C_r\). Even when each
individual rank gives a finite interval, approaching a zero leading
coefficient can make their union unbounded. A strict common margin is
what supplies the useful bound.

For illustration, take the previously studied error probability
\(e=.08\), so \(r=.84\), and allocate \(\alpha_H=.02\). At
\(N=100,q=.4\), the leading coefficient is \(-479/125000\); this
construction has unbounded tails even if the rank were known exactly. At
\(N=320,q=.4\), it is \(17009/2000000>0\) at that rank. With
\(N=100,q=.1\), the threshold exceeds one, so no allowable rank can
supply a positive leading coefficient. Exact rational arithmetic
reproduces these facts. This is a limitation of the conservative
Chebyshev construction, not an impossibility theorem for every sharper
mixture or randomization method. A bounded interval obtained by clipping
its tails would abandon the proved coverage statement.

Calendar uncertainty has a different dimension. With \(\rho=.95\),
stationary calendar standard deviation .6, and \(T=40\), the standard
deviation of the calendar average is about .4494; at \(T=160\) it is
.2776. Increasing \(N\) does not change either number. An assumed
persistence cap may improve a bound, and genuinely independent calendars
may supply new information. Neither benefit follows from creating more
account-date rows within the same calendar.

\hypertarget{an-information-counterexample-for-unrestricted-persistence}{%
\section{An information counterexample for unrestricted
persistence}\label{an-information-counterexample-for-unrestricted-persistence}}

Consider the favorable submodel with exact audits, no pair or cell
noise, and \(\beta=\delta=b\). Every signed focal and companion history
equals the same \(b+C_t\). Thus even infinitely many observed pairs
leave only one Gaussian calendar path. Let the calendar variance
\(v_c=\sigma^2>0\) and correlation \(\rho\) be known in this submodel,
making inference easier than in the intended unknown-nuisance problem.
For \(T\geq2\), the precision matrix of a stationary AR(1) correlation
matrix is tridiagonal. Summing its entries gives the mean information

\[
 I_T=\mathbf1^\top\Sigma_T^{-1}\mathbf1
   =\frac{T(1-\rho)+2\rho}{\sigma^2(1+\rho)}.
\]

The matrix identity is independently checked by exact rational inversion
for small \(T\), and the general formula follows by summing its two
endpoint diagonals, \(T-2\) interior diagonals and \(2(T-1)\)
off-diagonal entries. For fixed \(|\rho|<1\), \(I_T\) grows with \(T\).
But along the admissible stationary sequence \(\rho_T=1-T^{-2}\), it
converges to \(\sigma^{-2}\). Thus no increasing information about the
common mean is guaranteed uniformly over all persistence values below
one.

The limitation is not confined to unbiased estimators. Compare the laws
with means zero and \(\Delta\), the same known covariance, and
\(\Delta=\sigma\). Their total variation distance is

\[
 \operatorname{TV}(P_0,P_\Delta)
 =2\Phi\left(\frac{|\Delta|\sqrt{I_T}}2\right)-1.
\]

To see this, whiten the two Gaussian laws and project onto their
mean-difference direction. The likelihood-ratio threshold lies halfway
between the two projected means, leaving the stated difference of normal
tail probabilities. Along \(\rho_T\), the distance tends to
\(2\Phi(1/2)-1\simeq.3829\).

For any confidence set \(S_T\) with coverage at least \(1-\alpha\) at
both means, the definition of total variation gives

\[
 P_0(\Delta\in S_T)\geq1-\alpha-\operatorname{TV}(P_0,P_\Delta).
\]

Combining this with \(P_0(0\in S_T)\geq1-\alpha\) shows

\[
 P_0\{\operatorname{diam}(S_T)\geq|\Delta|\}
 \geq1-2\alpha-\operatorname{TV}(P_0,P_\Delta).
\]

At 95\% coverage the lower bound tends to about .5171. An honest
confidence set therefore cannot have uniformly vanishing diameter over
this persistence sequence, even with known covariance and exact
identity. This counterexample does not contradict pointwise asymptotic
normality for a fixed correlation, and it does not prove impossibility
at the previously studied specific fixed value .95. It rules out a
stronger uniform precision promise that lacks a persistence margin or
independent calendar information. In the monetary target, the diameter
scales by 1000 with no additional information gained.

\hypertarget{what-additional-information-would-change-the-decision}{%
\section{What additional information would change the
decision}\label{what-additional-information-would-change-the-decision}}

The finite region separates valid uncertainty from useful precision.
Bounds on unknown variances and persistence supply a model under which
the probability statement is established, but the resulting set can
still be unbounded. For a generated prototype those bounds may be
stipulated and tested as assumptions. For a later bank claim their
scientific origin must be established, or their estimation must receive
its own valid uncertainty budget. A successful calculation cannot
certify them.

The earlier 5,632-request study of the additive pair/calendar/gold law
found finite undercoverage in scarce-information and persistent-calendar
regimes despite favorable balanced-case results. Those fixed findings
motivate this alternative construction; they do not validate it. The
same study's wrong-inclusion, shared-audit and endogenous-label cases
remain outside the law used here. A valid set under independent
symmetric audits cannot be transported to those mechanisms by changing a
covariance estimator. \citep{exec13_persistent_predecessor}

There are three distinct next routes. A conservative region can serve as
a transparent reference in a generated prototype, with exact quadratic
and set-projection checks. A sharper mixture-likelihood or
test-inversion method might retain validity with less conservative
bounds, but it must account for the unknown variances and dependence
parameters. A changed information design can add independently audited
pairs, independent representative gold labels or genuinely independent
calendar clusters. Each addresses a different part of the uncertainty.
More gold labels alone cannot remove a shared-calendar obstruction,
while more dates alone do not make a persistent audit independent.

A prototype should preserve unbounded and disconnected results, measure
the cost of certified projection, and compare these information choices
under the same scientific target. The comparison can initially use
generated mechanisms without costly bank data. A later bank experiment
is justified when its information could narrow a decision-relevant
ambiguity enough to change an action. The present argument supports that
conditional choice; it does not promise a universally precise funding
estimate. The next servicing argument confronts a different information cost: learning action and observation probabilities while preserving the original eligible population. Both constructions budget uncertainty from information that a large row count alone cannot supply. \label{exec13:persistent-region-end}


% SOURCE docs/release/2026-10-29/source/docs/writing-repair/2026-10-29/src/0086cf696f3dce80.tex CHAR 0
\hypertarget{servicing-policies-with-estimated-action-and-observation-probabilities}{%

